\chapter{Vertigo}

\section*{Definition}
Vertigo is the feeling that you or the world around you are spinning, even when you are completely still. It is caused by a problem with your inner ear or brain.

\section*{What Happens in Your Body}
Vertigo can be peripheral (from the inner ear) such as BPPV, labyrinthitis, or Meniere disease, or central (from the brain) such as stroke or multiple sclerosis. Neurological symptoms like double vision or trouble speaking suggest a central cause.

\section*{Causes}
\begin{itemize}
  \item BPPV (most common)
  \item Vestibular neuritis
  \item Meniere disease
  \item Stroke/TIA
  \item Multiple sclerosis
\end{itemize}

\section*{When to See a Doctor}
See your doctor if:
\begin{itemize}
  \item Symptoms are severe or getting worse over time
  \item First episode or with neurologic deficits warrants immediate — seek medical evaluation
  \item You have other concerning symptoms such as fever, weight loss, or bleeding
\end{itemize}

\section*{Related Symptoms}
\begin{itemize}
  \item Nausea/vomiting
  \item Hearing loss/tinnitus
  \item Nystagmus
  \item loss of coordination
  \item Double vision
\end{itemize}
\textbf{Important:} If this symptom persists, your doctor will check for serious underlying causes.

\section*{Self-Care / Home Management}
\begin{itemize}
  \item During a vertigo episode, sit or lie still with your eyes closed in a safe position.
  \item Avoid sudden head movements, bending down, or looking up.
  \item Do not drive or operate machinery until the spinning sensation resolves completely.
  \item If you have been diagnosed with BPPV, your doctor may show you the Epley maneuver to perform at home.
  \item Reduce salt intake if you have Meniere's disease, as it can reduce vertigo attacks.
  \item Use a nightlight and remove tripping hazards from walkways to prevent falls.
\end{itemize}

\section*{How Your Doctor Will Evaluate You}
\begin{itemize}
  \item Detailed neurological history includes onset pattern, progression, triggers, associated symptoms, and functional impact.
  \item Comprehensive neurological examination assesses mental status, cranial nerves, motor/sensory function, reflexes, coordination, and gait.
  \item Neuroimaging (CT, MRI) is often indicated to evaluate for structural, vascular, or demyelinating causes.
  \item Specialized testing (EEG for seizures, nerve conduction studies for neuropathy, lumbar puncture for meningitis) is guided by clinical suspicion.
\end{itemize}

\section*{Treatment Options}
\begin{itemize}
  \item Treatment is directed at the underlying cause identified through diagnostic evaluation.
  \item Symptomatic pharmacotherapy may include analgesics, anticonvulsants, antidepressants, or disease-modifying agents as appropriate.
  \item Physical, occupational, and speech therapy play important roles in functional recovery and rehabilitation.
  \item Referral to neurology is indicated for unexplained, progressive, or treatment-refractory neurological symptoms.
\end{itemize}

\clearpage
